\documentclass[10pt,a4paper]{amsart}
\usepackage[margin=19mm]{geometry}
\usepackage{amsmath,amssymb,mathtools}
\usepackage[scr=rsfs,cal=euler]{mathalfa}
\usepackage{xfrac}
\usepackage[unicode,hidelinks,bookmarks=false]{hyperref}
\newcommand{\ie}{i.e.\ }
\newcommand{\cf}{cf.\ }
\newcommand{\resp}{respectively\ }
\newcommand{\Subsection}{\subsection}
\providecommand{\nobreakdash}{\nobreak\hbox{-}}
\setlength{\emergencystretch}{2em}
\multlinegap0pt
\usepackage{bm}
\usepackage{bbm}
\usepackage{dsfont}
\usepackage{xspace}
\usepackage{cancel}
\usepackage{mathtools}
\usepackage[shortlabels]{enumitem}
\usepackage{amsthm}
\makeatletter
\@ifundefined{claim}{\newtheorem{claim}{}}{}
\@ifundefined{theorem}{\newtheorem{theorem}{Theorem}}{}
\@ifundefined{lemma}{\newtheorem{lemma}[theorem]{Lemma}}{}
\makeatother
\makeatletter
\newcommand{\br}[1]{\left(#1\right)}
\newcommand{\cF}{\mathcal{F}}
\newcommand{\cI}{\mathcal{I}}
\newcommand{\cX}{\mathcal{X}}
\DeclareMathOperator{\cl}{cl}
\newcommand{\del}{\!\setminus\!}
\newcommand{\es}{\varnothing}
\renewcommand{\ss}{\subseteq}
\@ifundefined{subproof}{\newenvironment{subproof}[1][\proofname]{%
  \renewcommand{\qedsymbol}{$\blacksquare$}%
  \begin{proof}[Subproof:]%
}{%
  \end{proof}%
}}{}
\makeatother
\title[Modularity, Extensions and Connectivity in Infinite Matroids]{Modularity, Extensions and Connectivity in Infinite Matroids}
\author{Mattias Ehatamm, Peter Nelson, Fernanda Rivera Omana}
\date{Source: arXiv:2604.20778v1 (CC BY 4.0)}
\begin{document}
\maketitle

\noindent\emph{Source and licence.} Modularity, Extensions and Connectivity in Infinite Matroids. Version arXiv:2604.20778v1 (CC BY 4.0), distributed under the Creative Commons Attribution 4.0 International licence, as recorded in the arXiv licence metadata for this exact version and shown on the abstract page https://arxiv.org/abs/2604.20778v1. Full rights notice: licenses/arxiv-2604.20778-CC-BY-4.0.txt.\par\smallskip

\noindent\textit{Editorial scope.} Counted unit: the Theorem whose source label is \texttt{modcutextension} in the article (source0.tex lines 1829--1832, source label \texttt{modcutextension}), delivered together with the article's own complete original proof of it (source0.tex lines 1833--2041, one \texttt{proof} environment of 209 lines). No mathematics is written, repaired or re-derived: statement and proof are the article's own text line for line. Required context retained uncounted: modcutinter (lines 1784--1787). Every definition, display and label of those lines is retained unchanged. Omitted from this package and not redistributed: 1--1783, 1788--1828, 2042--4989. Nothing inside a retained excerpt is omitted or altered: every retained body is delivered exactly as the article's lines are written, so no omission receipt of any kind accompanies this package. Citations: the retained excerpts carry no citation command at all, so no bibliography entry of the article is transcribed and no reference is redistributed. Reference adaptation: none was needed and none was made; every reference inside the delivered unit resolves inside it, so no unresolved reference or citation is produced. Printed slips kept literal: no printed word, symbol, hypothesis or number of the counted statement or of its proof was altered. Layout and class: the article's own class is \texttt{amsart} with packages including amsmath, amssymb, enumerate, mathtools, xcolor, hyperref, fontenc, inputenc, lmodern, minted, newunicodechar; the wrapper is the lane's pinned \texttt{amsart} editorial wrapper at 10pt on A4, which restates the theorem-like environments the retained text needs and adds the article's own macro definitions that the retained text needs (\texttt{\textbackslash br}, \texttt{\textbackslash cF}, \texttt{\textbackslash cI}, \texttt{\textbackslash cX}, \texttt{\textbackslash cl}, \texttt{\textbackslash del}, \texttt{\textbackslash es}, \texttt{\textbackslash ss}), with extra wrapper packages (bm, bbm, dsfont, xspace, cancel, mathtools, enumitem). The delivered numbering is the unit's own. Assets: no retained span contains a figure environment, a graphics-inclusion command, a PDF-page inclusion, a table or a TikZ picture; no image file is copied into this package, the package declares no asset dependency and no image file is redistributed with it. Licence: version arXiv:2604.20778v1 of this article is distributed under the Creative Commons Attribution 4.0 International licence, as recorded in the arXiv licence metadata for this exact version and shown on the abstract page https://arxiv.org/abs/2604.20778v1. A full rights notice is delivered at licenses/arxiv-2604.20778-CC-BY-4.0.txt. Characters: every retained excerpt is pure ASCII, so the delivered engine, XeLaTeX with its default Latin Modern text font, reports no missing character.
\begin{lemma}\label{modcutinter}
  If $\cF$ is a modular cut of a matroid $M$, and $\cF_0 \subseteq \cF$ is a nonempty 
  modular family, then $\cap \cF_0 \in \cF$. 
\end{lemma}

\begin{theorem}\label{modcutextension}
  If $\cF$ is a modular cut in a matroid $M$, and $e \notin E(M)$,
  then there is a unique extension $M'$ of $M$ by $e$ for which $\cF_M^{M'} = \cF$. 
\end{theorem}

\begin{proof}
  Let $E = E(M)$, and let $\cX$ be the collection of subsets of $E$ whose closure is in $\cF$. 
  Note that any superset of an element of $\cX$ is in $\cX$. 
  Let $\cI_M$ be the collection of $M$-independent sets, 
  let $\cI_0 = \cI_M \cap \cX$, let $\cI_1 = \cI_M - \cX$, and let $\cI_e = \{I \cup \{e\} : I \in \cI_1\}$. 
  Finally, let $\cI = \cI_0 \cup \cI_1 \cup \cI_e$. 
  Our first claim reduces the problem to showing that $\cI$ is the collection of independent 
  sets of a matroid. 
  
  \begin{claim}\label{ismatroid}
    If $\cI$ is the collection of independent sets of a matroid $M'$ on $E \cup \{e\}$, 
    then $M'$ is the unique matroid for which $M' \del e = M$ and $\cF_M^{M'} = \cF$.
  \end{claim}
  \begin{subproof}
    Suppose that the given $M' = (E \cup \{e\}, \cI)$ is a matroid. 
    Since $\cI_M$ is the collection of sets in $\cI$ not containing $e$, we know that $M' \del e = M$. 

    Let $I$ be an independent set of $M = M' \del e$. 
    By construction, we have $\cl_M(I) \in \cF_{M}^{M'}$ if and only if $e \in \cl_{M'}(I)$, 
    which holds if and only if $I \cup \{e\} \notin \cI_e$, or in other words $I \in \cI_0$. 
    This is equivalent to the condition that $\cl_M(I) \in \cF$. Since this holds for every 
    independent set $I$, we can conclude that $\cF$ and $\cF_{M}^{M'}$ contain exactly 
    the same closures of independent sets, so are equal. 

    Suppose now that $M_1$ is any matroid with $M_1 \del e = M$ and $\cF_{M}^{M_1} = \cF$. 
    For each set $I \subseteq E$, the set $I \cup \{e\}$ is independent in $M_1$ 
    if and only if $e \notin \cl_{M_1}(I)$, which holds if and only if 
    $\cl_M(I) \notin \cF_{M}^{M_1} = \cF$. This implies that any two choices $M_1, M_1'$ for $M_1$ 
    have the same independent sets containing $e$ and satisfy $M_1 \del e = M_1' \del e$, 
    so are equal. Therefore the matroid $M'$ is unique with the stated properties. 
  \end{subproof}
  
  We show that $\cI$ is the indeed the collection of independent sets of a matroid $M'$ with ground set $E(M) \cup \{e\}$. 
  Note that $\cI \cap 2^{E} = \cI_{M} \cap 2^E$. 
  
  % Let $\cI_{\mathrm{max}}$ be the collection of maximal sets in $\cI$. 
  


  \begin{claim}\label{notcoloop}
    If $E \notin \cF$, then $(E \cup \{e\}, \cI)$ is a matroid. 
  \end{claim}
  \begin{subproof}
    Since $\cF$ is closed under taking superflats, we have $\cF = \es$, which implies
    that $\cI = \{I \ss 2^{E \cup \{e\}} : I - \{e\} \in \cI\}$. In other words, 
    $\cI$ is the collection of independent sets of the matroid obtained from $M$ by adding 
    $e$ as a coloop. 
  \end{subproof}


  By the last claim, we can now assume that $E \in \cF$. To verify the independence axioms for $\cI$ 
  is the collection of independent sets of a matroid, we need to characterize 
  when a set $I$ is maximal $\cI \cap 2^X$ for some $X$. The next two claims
  do this in the the case where $e \notin I$ and the case where $e \in I$. 


  \begin{claim}\label{foo1}
    Let $I \ss X \ss E(M) \cup \{e\}$ with $e \notin I$.
    Then $I$ is a maximal element of $\cI \cap 2^X$ if and only if 
    $I$ is a basis for $X - \{e\}$ in $M$, and if $e \in X$, then $X - \{e\} \in \cX$. 
  \end{claim}
  \begin{subproof}
    Suppose that $I$ is maximal in $\cI \cap 2^X$. Since $I \subseteq E(M)$ we have 
    $I \in \cI_M$ and $I \ss X - \{e\}$. If $I \subseteq J \subseteq X - \{e\}$ with $J \in \cI_M$, 
    then $J \in \cI \cap 2^X$, so $J = I$ by maximality. It follows that $I$ is a basis of $X - \{e\}$ in $M$. Furthermore, if $e \in X$ and $X - \{e\} \notin \cX$, so $I \notin \cX$, 
    which implies that $I \cup \{e\} \in \cI_e \cap 2^X$, contradicting the maximality of $I$. 
    So the forwards implication holds. 
    
    Conversely, suppose that $I$ is a basis for $X - \{e\}$ in $M$, and that if $e \in X$, then $X - \{e\} \in X$. Let $J \in \cI$ with $I \ss J \ss X$; we wish to show that $I = J$. 
    Since $J - \{e\} \in \cI_M$
    and $I \ss J - \{e\} \ss X - \{e\}$ and $I$ is a basis, we have $I = J - \{e\}$, 
    so it suffices to show that $e \notin J$. 
    
    If $e \in J$, then since $e \in X$ we have $X - \{e\} \in \cX$, 
    so $\cl_M(J - \{e\}) = \cl_M(I) = \cl_M(X - \{e\}) \in \cF$, and thus $J - \{e\} \in \cX$. 
    But the fact that $e \in J \in \cI$ implies that $J \in \cI_e$, so $J - \{e\} \notin \cX$, 
    a contradiction. 
  \end{subproof}

  \begin{claim}\label{foo2}
    Let $I \ss X \ss E \cup \{e\}$ with $e \in I$. Then $I$ is a maximal element of 
    $\cI \cap 2^X$ if and only if $I \in \cI_e$, and either 
    \begin{enumerate}[(i)]
      \item\label{xnotin} $X - \{e\} \notin \cX$ and $I - \{e\}$ is a basis of $X - \{e\}$ in $M$, or
      \item\label{xin} $X - \{e\} \in \cX$, and $\cl_M(I-\{e\})$ is a maximal proper subflat of $\cl_M(X - \{e\})$. 
    \end{enumerate}
  \end{claim}
  \begin{subproof}
    Suppose that $I$ is maximal in $\cI \cap 2^X$. Since $e \in I \in \cI$, 
    we have $I \in \cI_e$. 
    
    If $X - \{e\} \notin \cX$, then let $J \in \cI_M$ with
    $I - \{e\} \ss J \ss X - \{e\}$. Since $J \ss X - \{e\} \notin \cX$, 
    we have $J \notin \cX$, so $J \cup \{e\} \in \cI_e \cap 2^X \ss \cI \cap 2^X$, 
    so $I = J \cup \{e\}$ by the maximality of $I$. This implies that $J - \{e\} = I$; 
    it follows that $I$ is a basis of $X - \{e\}$ in $M$. 

    If $X - \{e\} \in \cX$, then let $I'$ be a basis of $X - \{e\}$ containing $I - \{e\}$. 
    For each $f \in I' - I$, since $I$ is maximal, we have $e \in I \cup \{f\} \notin \cI$, 
    which implies that $(I - \{e\}) \cup \{f\} \in \cX$. It follows that the modular cut 
    $\cF$ contains $\cl_M(I - \{e\} \cup \{f\})$ for all $f \in I' - I$. The set $I'$ 
    certifies that the flats $\cl_M(I - \{e\} \cup \{f\}) : f \in I' - I$ form a modular 
    family, and so their intersection lies in $\cF$ by Lemma~\ref{modcutinter}.
    This intersection has $I - \{e\}$ as a basis, which gives $I - \{e\} \notin \cX$, contradicting $I \in \cI_e$. 

    For the converse direction, suppose that (\ref{xnotin}) or (\ref{xin}) holds for $I$ and $X$. Let $J \in \cI$ with $I \ss J \ss X$; we wish to show that $J = I$, or equivalently 
    that $J - \{e\} = I - \{e\}$. Note that since $e \in J \in \cI$ we have $J \in \cI_e$. 
    
    If (\ref{xnotin}) holds, 
    Since $J \in \cI$ we have $J - \{e\} \in \cI_M$; 
    the fact that $I - \{e\}$ is a basis of $X - \{e\}$ gives that $J - \{e\} = I - \{e\}$ 
    and so $J = I$, as required. 

    Suppose now that (\ref{xin}) holds. For each $f \in J - I$, 
    since $F = \cl_M((I - \{e\} \cup \{f\}))$ is a proper superflat of $\cl_M(I - \{e\})$ 
    that is contained in $\cl_M(X - \{e\})$, it is equal to $\cl_M(X - \{e\})$ 
    and is thus in $\cF$. But this gives 
    $\cX \ni \cl_M((I - \{e\}) \cup \{f\}) \ss \cl_M(J - \{e\})$, 
    so $J - \{e\} \in \cX$, which contradicts $J \in \cI$. 
  \end{subproof}

  We now proceed to show that $\cI$ is the collection of independent sets of a matroid. 
  It is obvious that $\cI$ is nonempty and closed under taking subsets; 
  the next two claims verify the less trivial axioms. 

  \begin{claim}
    If $I$ is a non-maximal element of $\cI \cap 2^{E \cup \{e\}}$ and $B$ is a maximal element of 
    $\cI \cap 2^{E \cup \{e\}}$, then there exists $f \in B - I$ for which $I \cup \{f\} \in \cI$. 
  \end{claim}
  \begin{subproof}
    
    Suppose not. Note that $I$ is a maximal element of $\cI \cap 2^{I \cup B}$, 
    but there exists $f \in E - I \cup B$ such that $I \cup \{f\} \in \cI$. 

    % Since $B$ is maximal in $\cI \cap 2^{E \cup \{e\}}$, all three cases 
    % given by~\ref{foo1} and~\ref{foo2} imply that $B-\{e\}$ is independent in $M$
    % and that $B-\{e\}$ spans some hyperplane of $M$; let $H$ be a hyperplane 
    % of $M$ that is spanned by $B - \{e\}$. 
    
    \textbf{Case 1 : $e \notin I \cup B$}. By~\ref{foo1}, the set $I$ is a basis for $I \cup B$ in $M$
    and $B$ is a basis for $E$ in $M$, so $I \cup \{x\}$ is dependent in $M$, and so not in $\cI$,
    for each $x \in E-B$. Since $I \cup \{f\} \in \cI$ we have $f = e$, so $I \in \cI_1$ 
    and therefore $E = \cl_M(I) \notin \cF$, a contradiction. 

    \textbf{Case 2 : $e \in B - I$}. 
    By~\ref{foo1}, 
    the set $I$ is a basis for $I \cup (B - \{e\})$ and that $B - \{e\} \in \cX$. 
    But then $B \notin \cI_e$ and so $B \notin \cI$, contrary to the choice of $B$. 

    \textbf{Case 3: $e \in I - B$.}
    By~\ref{foo1}, the set $B$ is a basis for $E$ in $M$, so $\cl_M(B) = E \in \cF$ and thus $B \in \cX$ and so $(I \cup B) - \{e\} \in \cX$. Since $I$ is maximal in $\cI \cap 2^{I \cup B}$, 
    we get from ~\ref{foo2} that $\cl_M(I - \{e\})$ is a maximal proper subflat of $\cl_M((I \cup B)- \{e\})$. 
    
    The latter set contains $\cl_M(B) = E$, so $\cl_M(I - \{e\})$ is thus a hyperplane of $M$.
    Now $f \in E - (I \cup B)$ and $I \cup \{f\} \in \cI$, so $(I - \{e\}) \cup \{f\} \in \cI_1$. 
    But $I - \{e\}$ spans a hyperplane of $M$ and $(I - \{e\}) \cup \{f\}$ is independent, 
    so is a basis of $M$. Therefore $\cl_M(I - \{e\} \cup \{f\}) = E \in \cF$, 
    which contradicts $(I - \{e\}) \cup \{f\} \in \cI_1$. 

    \textbf{Case 4 : $e \in I \cap B$}. 
    The fact that $I,B \in \cI$ give that $I - \{e\} \notin \cX$ and $B - \{e\} \notin \cX$. 
    Since $I$ is a maximal element of $\cI \cap 2^{I \cup B}$, we get from~\ref{foo2} that 
    $I - \{e\}$ is a basis of $B - \{e\}$. Since $E \in \cX$, we also get from~\ref{foo2}
    that $\cl_M(B - \{e\})$ is a maximal proper subflat of $\cl_M(E) = E$, so 
    $\cl_M(B - \{e\})$ is a hyperplane of $M$.

    Now $f \notin I \cup B$ and $I \cup \{f\} \in \cI$, so $(I - \{e\}) \cup \{f\} \in \cI_1$
    and thus $f \notin \cl_M(I  - \{e\}) = \cl_M(B - \{e\})$. But $(I - \{e\})$ spans
    a hyperplane of $M$ and so $(I - \{e\}) \cup \{f\}$ is a basis of $M$. 
    Therefore $\cl_M((I - \{e\}) \cup \{f\}) = E \in \cF$, 
    contradicting $(I - \{e\}) \cup \{f\} \in \cI$. 
  \end{subproof}
  
  \begin{claim}
    If $X \ss E \cup \{e\}$, then each $I \in \cI \cap 2^X$ is contained in a maximal set in $\cI \cap 2^X$. 
  \end{claim}
  \begin{subproof}
    As before, we split into cases. 

    \textbf{Case 1 : $e \notin I$ and either $e \notin X$ or $X - \{e\} \in \cX$}.
    Since $e \notin \cI$, we have $I \in \cI_M$; let $J$ be a basis for $X$ in $M$ with $I \ss J$. 
    By~\ref{foo1}, the set $J$ is maximal in $\cI \cap 2^X$, as required. 

    \textbf{Case 2 : $e \in X - I$ and $X - \{e\} \notin \cX$.}
    Since $e \notin \cI$, we have $I \in \cI_M$; 
    let $J$ be a basis for $X - \{e\}$ in $M$ containing $I$. By~\ref{foo2}, the set $J \cup \{e\}$ 
    is maximal in $\cI \cap 2^X$, giving the result. 

    \textbf{Case 3 : $e \in I$ and $X - \{e\} \notin \cX$.}
    Since $e \in I \in \cI$, we have $I - \{e\} \in \cI_1$. Let $J$ be a basis for $X - \{e\}$ 
    in $M$ that contains $I - \{e\}$. 
    Since $J \ss X - \{e\} \in I_1$, we have $J \in \cI_1$ and so $J \cup \{e\} \in \cI_e$;
    By~\ref{foo2}, the set $J \cup \{e\}$ is 
    is maximal in $\cI \cap 2^X$, as required. 

    \textbf{Case 4 : $e \in I$ and $X - \{e\} \in \cX$.}
    Since $e \in I \in \cI$, we have $I - \{e\} \in \cI_1$. Let $J$ be a basis for $X - \{e\}$ 
    in $M$ that contains $I - \{e\}$. If there is some $f \in J$ for which 
    $\cl_M(J - \{f\}) \notin \cF$, then $J - \{f\} \in \cI_e$, 
    and so $J \cup \{e\}  - \{f\}$ contains $I$ and  is maximal in $\cI \cap 2^X$ by~\ref{foo2}. 

    Otherwise, we have $\cl_M(J \cup \{e\} - \{f\}) \in \cF$ for all $f \in J - (I - \{e\})$; 
    the collection of all such flats is modular as certified by the independent set $J$, 
    so using Lemma~\ref{modcutinter} we have
    \[\cF \ni \bigcap_{f \in J - (I - \{e\})} \cl_M(J - \{f\}) = \cl_M\br{\bigcap_f (J - \{f\})} = \cl_M(I - \{e\}).\]
    This contradicts $I \in \cI$. 
  \end{subproof}
  By the last two claims and~\ref{ismatroid}, the theorem holds. 
\end{proof}

\end{document}
